\documentclass[11pt]{article}


\usepackage{a4wide}
\usepackage{amssymb}
\usepackage[utf8]{inputenc}
\usepackage[russian]{babel}
\usepackage{graphicx}
\usepackage{subcaption}
\usepackage{amsmath}
\usepackage{mathtools}% Для улучшенной работы с математикой

\newcommand{\R}{\ensuremath{\mathbb{R}}}
\newcommand{\norm}[1]{\left\lVert #1 \right\rVert}
\newcommand{\scalar}[2]{\left<#1,#2\right>}
\newcommand{\opor}[2]{\rho(#1\; | \;#2)}
\newcommand{\supr}[2]{\sup\limits_{#1 \in #2}}
\renewcommand{\epsilon}{\varepsilon}
\renewcommand{\phi}{\varphi}
\renewcommand{\le}{\leqslant} 
\renewcommand{\ge}{\geqslant}
\newcommand{\R}{\mathbb{R}}
\newcommand{\Lcal}{\mathcal{L}}
\newcommand{\Rnn}{\mathbb{R}^{2 \times 2}}%R^(2*2)
\newcommand{\Pcal}{\mathcal{P}}%P -ограничение на u
\newcommand{\Xcal}{\mathcal{X}}
\begin{document}

\thispagestyle{empty}

\begin{center}
\ \vspace{-3cm}

\includegraphics[width=0.5\textwidth]{msu.eps}\\
{\scshape Московский государственный университет имени М.~В.~Ломоносова}\\
Факультет вычислительной математики и кибернетики\\
Кафедра системного анализа

\vfill

{\Huge\bfseries <<Стохастический анализ и моделирование>>}
\end{center}

\vspace{1cm}

\begin{flushright}
  \large
  \textit{Студент 415 группы}\\
  А.\,Д.~Тулегенов

  \vspace{5mm}

  \textit{Руководитель практикума}\\
  аспирант В.\,А.~Сливинский
\end{flushright}

\vfill

\begin{center}
Москва, 2025
\end{center}

\newpage
% Содержание
\tableofcontents
\newpage
%начало по делу
\section{Задание 1}
\subsection{Постановка задачи}
\begin{enumerate}
  \item Реализовать генератор схемы Бернулли с заданной вероятностью успеха $p$. На основе генератора схемы Бернулли построить датчик биномиального распределения.
  \item Реализовать генератор геометрического распределения; проверить для данного распределения свойство отсутствия памяти.
  \item Промоделировать игру в орлянку: бесконечную последовательность независимых испытаний Бернулли с бросанием “правильной” (честной, $p = 0.5$) монеты. Величина “выигрыша” $S_n$ определяется как сумма по $n$ испытаниям значений $1$ и $-1$ в зависимости от выпавшей стороны монеты. Проиллюстрировать в виде ломаной поведение нормированной суммы
  \[
    Y(i)=\frac{S_i}{\sqrt{i}}
  \]
  как функцию от номера испытания $i$ для отдельно взятой траектории. Дать теоретическую оценку для значения $Y(n)$ при $n\to\infty$.
\end{enumerate}
\subsection{Теоретические сведения}
\subsubsection{Схема Бернулли и биномиальное распределение}
\textbf{Определение.} Схемой Бернулли с заданной вероятностью $p$ называется эксперимент, состоящий из серии испытаний, удовлетворяющих следующим условиям:
\begin{enumerate}
  \item Отсутствие взаимного влияния в испытаниях;
  \item Воспроизводимость. Однородные испытания проводятся в сходных условиях;
  \item Существует признак, который реализуется(успех с вероятностью p) или не реализуется(неуспех с вероятностью 1-p) в испытании. Признак может быть отнесён к любому из испытаний.
\end{enumerate}
\[
\mathbf{P}(x = k) =
\begin{cases}
1 - p, & k = 0 \\
p, & k = 1
\end{cases}
\]
Под генератором и датчиком везде будем подразумевать функцию, возвращающую выборку заданного размера из соответствующего распределения.

\textbf{Определение.} Биномиальное распределение описывает число успехов в $n$ независимых испытаниях схемы Бернулли. Вероятность получить $k$ успехов определяется формулой:
\[
  \mathbb{P}(X=k)=\binom{n}{k}p^k(1-p)^{n-k},\quad k=0,1,\dots,n.
\]

\subsubsection{Геометрическое распределение}
Геометрическое распределение описывает число неудачных испытаний до первого успеха при независимых повторениях с постоянной вероятностью успеха.$\newline$
Функция вероятности геометрического распределения имеет вид $\mathbb{P}(X = n) = p(1 - p)^n. \newline$
Получим геометрическое распределение через экспоненциальное. Пусть $\xi \sim exp(\lambda)$, тогда $\lfloor \xi \rfloor \sim Geom(p)$, где $p = 1 - e^{-\lambda}$.
Покажем, что это так. Функция распределения имеет вид $F_\xi(x) = 1-e^{-\lambda x}$. Поэтому для $n \in \mathbb{N}_0$
$$
    \mathbb{P}(\lfloor \xi \rfloor = n) = \mathbb{P}(n \le \xi \le n + 1) = F_\xi(n+1) - F_\xi(n) =  e^{-n\lambda} (1 - e^{-\lambda}) = (1 - p)^np
$$
Таким образом, моделирование геометрического распределения с параметром p сводится к моделированию экспоненциального распределения с параметром $\lambda = - \ln(1 - p)$.
Воспользуемся методом обращения функции распределения:
$$
    F_\xi^{-1}(y) = -\frac{1}{\lambda}\ln(1 - y)
$$
Пусть случайная величина $\xi$ имеет непрерывную и строго возрастающую функцию распределения $F_\xi$. Тогда $F_\xi^{-1}(U) \sim \xi,\ U \sim \mathcal{U}(0, 1)$
$$
    \mathbb{P}(F_\xi^{-1}(U) < x) = \mathbb{P}(U < F_\xi(x)) = F_\xi(x),
$$
поэтому $F_\xi^{-1}(U) \sim \xi$

Теперь проверим, что для геометрического распределения выполняется свойство отсутствия памяти, т.е.: $$\mathbb{P}(X > m + n \mid X \geqslant n) = \mathbb{P}(X > m), \: X \sim Geom(p)$$


Мы имеем: $$\mathbb{P}(X > m) = (1 - p)^{m + 1}$$

Поэтому мы получим: $$\mathbb{P}(X > m + n \mid X \geqslant n) = \frac{\mathbb{P}(X > m + n, X \geqslant n)}{\mathbb{P}(X \geqslant n)} = \frac{(1 - p)^{m + n + 1}}{(1 - p)^n} = \mathbb{P}(X > m)$$
Что и требовалось доказать.
\subsubsection{Игра в орлянку}
Игра в орлянку - бесконечная последовательность независимых испытаний с бросанием правильной монеты. Выйгрыш $S_n$ определяется как сумма по n испытаниям значений 1 и -1 в значении от выпавшей стороны. Проиллюстрируем поведение величины $Y(i) = \frac{S_i}{\sqrt n}$ как функцию номера испытания $i = 1, \dots ,n$ для отдельно взятой траектории.\\
\textbf{Теорема (ЦПТ)}
При $n \to \infty$ нормированная сумма $Y(i)$ стремится к нормальному распределению с математическим ожиданием $0$ и дисперсией $1$:
$$
Y(i) \xrightarrow{d} N(0, 1).
$$
Согласно теореме, при большом числе испытаний ($i \to \infty$) $Y(i)$ приближается к стандартному нормальному распределению $N(0, 1)$. Моделировали эту нормированную сумму на практике, используя последовательность случайных испытаний Бернулли с вероятностью успеха $p = 0.5$.
\subsection{Результаты:}
\subsubsection{Схема Бернулли и биномиальное распределение}
\begin{center}
\includegraphics[width=0.7\linewidth]{мой отчёт/1.png}
\end{center}

\subsubsection{Проверка отсутствия памяти:}
\begin{center}
\includegraphics[width=0.7\linewidth]{мой отчёт/2.png}
\end{center}

\subsubsection{Игра в орлянку:}
\begin{center}
\includegraphics[width=0.7\linewidth]{мой отчёт/3.png}
\caption{График нормированной суммы $Y(i)=S_i/\sqrt{i}$.}
\end{center}
\section{Задание 2}
\subsection{Постановка задачи}
\begin{enumerate}
\item Построить датчик сингулярного распределения, имеющий в качестве функции распределения канторову лестницу. С помощью критерия Колмогорова убедиться в корректности работы датчика.
\item Для канторовых случайных величин с помощью критерия Смирнова проверить:
\end{enumerate}
\begin{itemize}
\item свойство симметричности относительно $\tfrac12$, то есть случайные величины $X$ и $1-X$ распределены одинаково;
\item свойство самоподобия относительно деления на 3, то есть условное распределение $Y$ при условии $Y\in[0,\tfrac13]$ совпадает с распределением $Y/3$.
\end{itemize}
\item Рассчитать значения математического ожидания и дисперсии для данного распределения. Сравнить теоретические значения с эмпирическими и проиллюстрировать сходимость.
\subsection{Теоретические сведения}
\subsubsection{Сингулярное распределение}
Построим датчик сингулярного распределения с функцией распределения, совпадающей с канторовой лестницей.

Из определения канторовой лестницы следует, что каждая цифра в троичной записи значения случайной величины принимает значения $0$ или $2$ с вероятностью $\frac{1}{2}$.

Введём последовательность независимых одинаково распределённых случайных величин
$$
\xi_i \in \{0, 2\}, \qquad \mathbb{P}(\xi_i = 0) = \mathbb{P}(\xi_i = 2) = \frac{1}{2}.
$$

Тогда случайную величину с канторовым распределением можно представить в виде
$$
X = \sum_{i=1}^{\infty} \frac{\xi_i}{3^i}.
$$

В численных экспериментах будем использовать конечную сумму
$$
X_{20} = \sum_{i=1}^{20} \frac{\xi_i}{3^i},
$$
что обеспечивает достаточную точность аппроксимации канторова распределения.\\
Функция Кантора имеет вид
$$
C(x) =  \sum_{i=1}^{\infty} \frac{\xi_i}{2^i}.
$$
\subsubsection{Проверка критерием Колмогорова}
Критерий Колмогорова:

Установим уровень значимости $\alpha$ = 0.05
1. Вычисляется максимальное отклонение:
$$
    D = \sup_x \lvert С_{emp}(x) - С(x) \rvert, где~С_{emp}(x)~и~С(x) - эмпирическая~и~истинная~функции~распределения.
$$
2. Используется статистика $ F_k $ для вычисления $ p $-value:
$$
    F_k = \sqrt(n) D = 1 + 2 \sum_{k=1}^\infty (-1)^{k} e^{-2k^2 D^2},    
$$
то есть:
$$
    F_k = 1 - 2 \sum_{k=1}^\infty (-1)^{k-1} e^{-2k^2 D^2},    
    p = 1 - F_k.
$$

Решение:\\
1. Если $ p > \alpha $, гипотеза принимается (расхождения несущественны).\\
2. Если $ p \leq \alpha $, гипотеза отвергается (расхождения значимы).


\subsubsection{Проверка критерием Смирнова}
Двухвыборочный критерий Смирнова (KS-test) сравнивает две выборки и измеряет
$$D_{n,m}=\sup_x\lvert \hat F_{1,n}(x)-\hat F_{2,m}(x)\rvert.$$
Если распределения одинаковы, то при больших $n,m$ статистика $\sqrt{\frac{nm}{n+m}}D_{n,m}$ имеет колмогоровское распределение.

Нам нужно проверить два свойства:\\
1) Симметрия относительно $1/2$: $X\overset{d}=1-X$.
2) Самоподобие относительно деления на 3: условное распределение $X\mid\{X\in[0,1/3]\}$ совпадает с распределением $X/3$.

Симметрия «относительно $1/2$» означает, что отражение точки $x$ относительно $1/2$ даёт $1-x$.
Поэтому достаточно проверить однородность распределений выборок из $X$ и из $1-X$.

Для самоподобия при делении на 3 сравним $X\mid\{X\le 1/3\}$ и $X/3$

Строим две выборки одинакового размера:\\
- $x_cond$: из условного распределения $X\mid\{X\le 1/3\}$.\\
- $x_scaled$: из распределения $X/3$ (берём независимую выборку $X$ и делим на 3).\\
Дальше, аналогичным образом пользуемся критерием Смирнова
\subsubsection{Нахождение матожидания и дисперсии}
## Математическое ожидание и Дисперсия
Найдём эмпирические значения матожидания и дисперсии.
Рассмотрим случайную величину $\xi_i \sim 2*\text{Bernoulli}(0.5)$, поэтому $\mathbb{E}[\xi_i] = 1$. Математическое ожидание $\mathbb{E}[X]$ выражается как:

$$
\mathbb{E}[X] = \sum_{i=1}^{\infty} \mathbb{E}[\xi_i] \cdot 3^{-i}.
$$

Сумма прогрессии:

$$
\sum_{i=1}^{\infty} 3^{-i} = \frac{1}{3-1} = \frac{1}{2}.
$$

Таким образом:

$$
\mathbb{E}[X] = 1 \cdot \frac{1}{2} = \frac{1}{2}.
$$
Аналогично, дисперсия $\text{Var}(X) = \frac{1}{8}$.

\subsection{Результаты:}
\subsubsection{Сингулярное распределение}
\begin{center}
\includegraphics[width=0.7\linewidth]{мой отчёт/4.png}
\caption{Истинная и эмпирическая функции распределения}
\end{center}

\subsection{Проверка критерием Колмогорова}
Установим уровень значимости $\alpha = 0.1$. Проведём 
300 испытаний при сгенерированной выборке размером $n = 1000$ и посмотрим в какой доле случаев полученное p-value меньше $\alpha$. В результате работы программы получаем, что гипотеза отклоняется в 4.1\% случаев, что является очень хорошим результатом. 
\subsection{Проверка критерием Смирнова}
Установим уровень значимости $\alpha = 0.1$. Аналогично, но для критерия Смирнова получим отклонение гипотезы в 4\% cлучаев.\\
Оба критерия дают очень хороший результат для данного распределения.
\subsection{Нахождение матожидания и дисперсии}
Сгенерируем выборку размером в 100000. Получим, что выборочное матожидание равно 0.5001182968026713, а выборочная дисперсия равна 0.12510173057982424. Проиллюстрируем сходимость эмпирического матожидания к теоретическому, и аналогично для дисперсии:
\begin{center}
\includegraphics[width=0.7\linewidth]{мой отчёт/5.png}
\end{center}

\begin{center}
\includegraphics[width=0.7\linewidth]{мой отчёт/6.png}
\caption{Выборочные матожидание и дисперсия}
\end{center}
По графику видно, что обе эмпирические величины хорошо сходятся к теоретическим.

\section{Задание 3}
\subsection{Постановка задачи}
\begin{enumerate}
\item Построить датчик экспоненциального распределения. Проверить для данного распределения свойство отсутствия памяти.
\item Пусть $X_1,\dots,X_n$ — независимые экспоненциально распределённые случайные величины с параметрами $\lambda_1,\dots,\lambda_n$. Найти распределение $Y=\min(X_1,\dots,X_n)$.
\item На основе датчика экспоненциального распределения построить датчик пуассоновского распределения.
\item Построить датчик пуассоновского распределения как предел биномиального. Убедиться в корректности построенного датчика при помощи критерия $\chi^2$ Пирсона.
\item Построить датчик стандартного нормального распределения методом моделирования случайных величин парами с переходом в полярные координаты (преобразование Бокса–Мюллера). Проверить при помощи t-критерия Стьюдента равенство математических ожиданий, а при помощи критерия Фишера — равенство дисперсий.
\end{enumerate}
\subsection{Теоретические сведения}
\subsubsection{Экспоненциальное распределение}
Датчик экспоненциального распределения был получен в задаче 1 методом обращения функции распределения. Свойство отсутствия последействия формулируется и доказывается в этом случае аналогично тому, как это было сделано для геометрического
распределения.
\subsubsection{Распределение $Y = \min(X_1, \dots, X_n)$}
\begin{equation*}
\begin{aligned}
F_Y(x) = \mathbb{P}(\min(X_1, \dots, X_n) < x) &= 1 - \mathbb{P}(X_1 \ge x, \dots, X_n \ge x) = \\
&= 1 - \prod_{i=1}^n (1 - F_{X_i}(x)) = 1 - \exp \left( -\sum_{i=1}^n \lambda_i x \right),
\end{aligned}
\end{equation*}
откуда $Y \sim \exp \left( \sum_{k=1}^n \lambda_k \right)$.
\subsubsection{Датчик пуассоновского распределения через экспоненциальное}
Распределение количества событий N(t) в пуассоновском процессе:
$$
P(N(t) = k) = \frac{(\lambda t)^k}{k!} e^{-\lambda t}, \quad k = 0, 1, 2, \dots
$$
При t = 1 формула упрощается и совпадает с классическим распределением Пуассона:
$$
P(N(1) = k) = \frac{\lambda^k}{k!} e^{-\lambda}
$$
То есть, достаточно посчитать количество событий за единицу времени.
\subsubsection{Датчик пуассоновского распределения через биномиальное}
Распределение биномиальной случайной величины:
$$
P(X = k) = C_n^k p^k (1 - p)^{n-k}, \quad k = 0, 1, 2, \dots, n,
$$

Биномиальное распределение при $ n \to \infty, \, p \to 0, \, n \cdot p = \lambda $:
$$
P(X = k) \approx \frac{\lambda^k e^{-\lambda}}{k!}, \quad \lambda = n \cdot p.
$$

Проверим правильность при помощи критерия согласия Пирсона. В нем используется статистика хи-квадрат:
$$
\chi^2 = n \sum_{i = 0}^k \frac{(\frac{n_i}{n} - \mathbb{P}_i(X))^2}{\mathbb{P}_i(X)},
$$
где $n_i$ - число значений, равных $i$, $\mathbb{P}_i(X) = \mathbb{P}(X = i)$ при $i = \overline{0,k-1}$; $n_k$ - число значений, больших либо равных $k$, $\mathbb{P}_k(X) = \mathbb{P}(X \ge k)$.

Зафиксируем уровень значимости $\alpha = 0.05$. Нулевая гипотеза о том, что выборка порождена пуассоновским распределением принимается, если значение функции распределения $\chi^2_{k-1}$ на статистике $\chi_2$ больше уровня значимости, иначе она отвергается. k выбирается так, чтобы при каждом значении $i < k$ было достаточно элементов выборки, равных $i$ (>5).
\subsubsection{Датчик стандартного распределения}
Если $U_1, U_2 \sim U(0,1)$ --- независимые равномерно распределённые случайные величины, то стандартное нормальное распределение $Z_1, Z_2 \sim N(0,1)$ можно получить по формулам:
\[
Z_1 = \sqrt{-2 \ln U_1} \cos(2\pi U_2), \quad Z_2 = \sqrt{-2 \ln U_1} \sin(2\pi U_2).
\]
Алгоритм метода Бокса-Мюллера состоит в следующем:

\begin{enumerate}
    \item Сгенерировать две независимые случайные величины $U_1, U_2 \sim U(0,1)$.
    \item Вычислить радиус $R$ и угол $\theta$ в полярных координатах:
    \[
    R = \sqrt{-2 \ln U_1}, \quad \theta = 2\pi U_2.
    \]
    \item Перейти в декартовы координаты:
    \[
    Z_1 = R \cos(\theta), \quad Z_2 = R \sin(\theta).
    \]
\end{enumerate}

Полученные $Z_1$ и $Z_2$ являются независимыми стандартно-нормальными случайными величинами.\\
\textbf{Критерий t-Стьюдента}\\
Критерий Стьюдента используется для проверки гипотезы о равенстве мат. ожидания выборки заданному значению:
$$
H_0: \mu = \mu_0, \quad H_1: \mu \neq \mu_0.
$$
Статистика t-критерия
$$
t = \sqrt n \frac{\bar{X} - \mu_0}{S}
$$
где:
- $\bar{X} = \frac{1}{n} \sum_{i=1}^n X_i$ - выборочное среднее
- $S = \sqrt {\frac{1}{n-1} \sum_{i = 0}(X_i - \bar{X})}$ - стандартное несмещённое отклонение
- $n$ - объём выборки\\

\textbf{Критерий Фишера}\\
Критерий Фишера используется для проверки гипотезы о равенстве
дисперсий двух выборок:
$$
H_0:S_1^2 = S_2^2, \quad H_1:S_1^2 \neq S_2^2
$$
Статистика:
$$
F = \frac{S_1^2}{S_2^2}
$$
Критерий Фишера является двусторонним.
\subsection{Результаты}

\begin{center}
\includegraphics[width=0.7\linewidth]{мой отчёт/7.png}
\caption{Моделирование экспоненциального распределения при размере выборки $N = 1000$.}
\end{center}

\subsubsection{Проверка отсутствия памяти}

\begin{center}
\includegraphics[width=0.7\linewidth]{мой отчёт/8.png}
\caption{Моделирование экспоненциального распределения при размере выборки $N = 1000$.}
\end{center}

\subsubsection{Распределение $Y = \min(X_1, \dots, X_n)$}

\begin{center}
\includegraphics[width=0.7\linewidth]{мой отчёт/9.png}
\caption{Моделирование экспоненциального распределения при размере выборки $N = 1000$.}
\end{center}

Установим уровень значимости $\alpha = 0.05$.\\
Применим критерий согласия Пирсона при выборке в $n = 1000$ и кол-ве испытаний в 1000. Гипотеза $H_0$ отвергается 71 раз из 1000, то есть гипотеза верна в 93\% случаях.

\subsubsection{Стандартное нормальное распределение}

\begin{center}
\includegraphics[width=0.7\linewidth]{мой отчёт/10.png}
\caption{Моделирование экспоненциального распределения при размере выборки $N = 1000$.}
\end{center}

При тех же значениях размерах выборки и уровне значимости $\alpha$ проверим выполнение ещё двух статистических критерия.\\

С помощью t-критерия Стьюдента проверим равенство матожидания выборки теоретическому, то есть нулю. Гипотеза $H_0$ отвергается 57 раз из 1000, то есть гипотеза верна в 95\% случаях. Таким образом, нулевую гипотезу можно считать верной.\\


С помощью критерия Фишера проверим равенство дисперсий двух выборо теоретическому, то есть нулю. Гипотеза $H_0$ отвергается 47 раз из 1000, то есть гипотеза верна в 95\% случаях. Таким образом, нулевую гипотезу можно считать верной.

\section{Задание 4}
\subsection{Постановка задачи}

\begin{enumerate}
\item Построить датчик распределения Коши.
\item На основе датчика распределения Коши с помощью метода фон Неймана построить датчик стандартного нормального распределения. При помощи графика normal probability plot убедиться в корректности построенного датчика и обосновать наблюдаемую линейную зависимость.
\item Сравнить скорость моделирования стандартного нормального распределения в заданиях 3 и 4.
\end{enumerate}

\subsection{Теоретические сведения:}
\subsubsection{Датчик распределения Коши}
Функция Распределения случайной величины, имеющей распределение Коши $C(x_0, \gamma)$ с параметрами $x_0 \in \mathbb R$ и $\gamma > 0$ имеет следующую функцию распределения:
$$
    F_X(x) = \frac{1}{\pi}\arctan(\frac{x - x_0}{\gamma}) + \frac{1}{2}
$$
где:
- $x_0 \in \mathbb R$ - параметр сдвига
- $\gamma > 0$ - параметром масштаба

Воспользуемся методом обратной функции распределения.

Обратная функция распределения Коши имеет вид:
$$
F^{-1}(y) = x_0 + \gamma \tan\!\left( \pi \left( y - \frac{1}{2} \right) \right).
$$

Будем использовать её для моделирования распределения Коши:
$$
\xi = F^{-1}(U) \sim \mathcal{C}(x_0, \gamma),
\qquad U \sim Uniform(0,1).
$$
\subsubsection{метод Фон Неймана}

Метод фон Неймана (accept–reject, rejection sampling) используется для моделирования случайных величин $X \sim \mathcal N(0,1)$ с заданной плотностью $f(x)$ в ситуации, когда прямое моделирование затруднено, но доступен датчик для другого распределения с плотностью $g(x)$.$\newline$
Идея метода состоит в следующем. Пусть требуется получить выборку из распределения с плотностью $f(x)$, а мы умеем генерировать случайные величины с плотностью $g(x)$. Предполагается, что существует такая константа $M > 0$, что
$$
f(x) \le M g(x), \quad \forall x \in \mathbb{R}.
$$
Для каждой случайной выборки с плотностью $f(x)$ мы генерируем значение случайной величины и принимаем его в случае, если оно удовлетворяет неравенству:
$$
U \leq \frac{f(Y)}{M g(Y)}
$$
где
- $ U \sim \text{Uniform}(0, 1) $
- $Y \sim \mathcal C(0,1)$.

Величину $Y$ моделируем методом обратной функции:
$$
Y=\tan\bigl(\pi(U-\tfrac12)\bigr), \quad U\sim Uniform(0,1).
$$
Теоретическое распределение Коши:
$$
f_C(x) = \frac{1}{\pi(1 + x^2)}
$$
Нормальное распределение:
$$
f_N(x) = \frac{e^\frac{-x^2}{2}}{\sqrt(2\pi)}
$$

Будем искать минимальную константу $M$, так как это влияет на скорость метода. Будем искать супремум отношения $r(x) = \frac{f(x)}{g(x)}$.
$$
r(x) = \frac{f(x)}{g(x)} = \sqrt\frac{\pi}{2}(x^2 + 1)e^{-\frac{x^2}{2}}
$$
$$
r'(x) = \sqrt\frac{\pi}{2}\left(xe^{\frac{-x^2}{2}} - (x^2 + 1)\frac{x^2}{2}xe^{-\frac{x^2}{2}}\right) = \sqrt\frac{\pi}{2}xe^{-\frac{x^2}{2}}\frac{1}{2}\left(2 - x^4 - x^2\right) = 0 \Rightarrow x = 0,\, x = \pm 1.
$$
$$
r(0) = \sqrt\frac{\pi}{2},\, \quad r(\pm 1) = \sqrt\frac{2\pi}{e}.
$$
Так как $M = sup_x r(x)$, то из:
$$
 \sqrt\frac{2\pi}{e} > \sqrt\frac{\pi}{2} \Rightarrow M = \sqrt\frac{2\pi}{e}.
 $$

 Смысл метода в том, что мы генерируем значение из известного распределения $x$ и принимаем его в выборку целевого распределения с вероятностью $\frac{f(x)}{Mg(x)}$. Обоснуем применимость метода. Рассмотрим плотность вероятности при условии, что мы приняли значение:
$$
    p(x|A) = \frac{\mathbb{P}(A|x)p(x)}{\mathbb{P}(A)} = \frac{\frac{f(x)}{Mg(x)}g(x)}{\mathbb{P}(A)}
$$
$$
    \mathbb P(A) = \int_{-\infty}^{+\infty} \frac{f(x)}{Mg(x)} g(x) \mathrm dx = \int_{-\infty}^{+\infty} \frac{f(x)}{M} \mathrm dx = \frac{1}{M}
$$
$$
    p(x|A) = \frac{\frac{f(x)}{Mg(x)}g(x)}{\frac{1}{M}} = f(x)
$$
Таким образом, мы получаем нужную плотность распределения.
\subsection{normal probability plot}\\
Normal Probability Plot это график, на котором отображаются квантили выборки (по оси $y$) против квантилей теоретического нормального распределения (по оси $x$). Если выборка распределена нормально, точки на графике будут лежать на прямой линии.$\newline$
Любое нестандартное нормальное распределение можно представить в виде $\mu + \sigma \mathcal N(0, 1)$, то есть в виде линейного преобразования стандартного, поэтому график также будет линейным. Для других распределений квантили ведут себя по другому, поэтому график не будет линейным. Также графики не будут инвариантны относительно параметров масштаба и сдвига.

\subsection{Результаты:}
\subsubsection{Распределение Коши}
\begin{center}
\includegraphics[width=0.7\linewidth]{мой отчёт/11.png}
\end{center}

\subsubsection{Стандартное нормальное распредлеение через метод фон Неймана}
\begin{center}
\includegraphics[width=0.7\linewidth]{мой отчёт/12.png}
\end{center}

\begin{center}
\includegraphics[width=0.7\linewidth]{мой отчёт/13.png}
\end{center}

\subsubsection{Проверка корректности с помощью normal probability plot}
\begin{center}
\includegraphics[width=0.7\linewidth]{мой отчёт/14.png}
\end{center}
На графике можно наблюдать как точки хорошо ложаться на прямую, что подтверждает корректность построенного датчика стандартного нормального распределения.
\begin{center}
\includegraphics[width=0.7\linewidth]{мой отчёт/15.png}
\end{center}
По графику видно, что метод Бокс-Мюллера работает значительно быстрее.

\section{Задание 5}
\subsection{Постановка задачи}
\begin{enumerate}
  \item Пусть $X_i \sim N(\mu, \sigma^2).$ Убедиться эмпирически в справедливости теоремы о законе больших чисел (ЗБЧ) и центральной предельной теоремы (ЦПТ): исследовать поведение суммы $S_n = \sum_{i=1}^n X_i$ и эмпирического распределения величины
$$
\sqrt{n} ( \frac{S_n}{n} - a). 
$$
  \item Считая $\mu$ и $\sigma^2$ неизвестными, построить доверительные интервалы для среднего и дисперсии по имеющейся выборке.
  \item Пусть $X_i \sim K(a,b)$ — имеет распределение Коши с параметром сдвига $a$ и масштаба $b$. Изучить эмпирически как ведут себя суммы $S_n/n$, объяснить результат и найти закон распределения данных сумм.
\end{enumerate}
\subsection{Теоретические сведения}
\subsubsection{ЗБЧ и ЦПТ}
Рассмотрим независимые одинакого распределенные случайные величины: $X_i \sim \mathcal{N} (\mu, \sigma^2)$. Убедимся эмприрически в справедливости закона больших чисел (ЗБЧ). Рассмотрим выборочное среднее $\bar{X}_n = \frac{X_1 + \dots + X_n}{n}$, и покажем что $\bar{X}_n \xrightarrow[n \to \infty]{p} \mu$.\\
ЦПТ, гласит что $Y = \frac{S_n - \mu n}{\sigma \sqrt n} \xrightarrow[n \to \infty]{d} \mathcal{N}(0, 1)$, где $S_n = X_1 + \dots + X_n$. Сходимость по распределению означает, что распределение получившейся величины близко к целевому.
\subsubsection{Доверительные интервалы для матожидания и дисперсии}
Пусть $X_1,\dots,X_n$ — независимые случайные величины, $X_i \sim \mathcal N(\mu,\sigma^2)$, где параметры $\mu$ и $\sigma$ неизвестны.

Обозначим выборочные оценки:
$$
\bar X = \frac{1}{n}\sum_{i=1}^n X_i,\qquad
S^2 = \frac{1}{n-1}\sum_{i=1}^n (X_i-\bar X)^2.
$$

\textbf{Доверительный интервал для $\mu$ (уровень доверия $1-\alpha$)}

Для нормальной выборки при неизвестной $\sigma$ верно:
$$
T=\frac{\bar X-\mu}{S/\sqrt{n}} \sim t_{n-1}.
$$
Отсюда доверительный интервал:
$$
\mu \in \left[\bar X - t_{1-\alpha/2,n-1}\frac{S}{\sqrt{n}},\;
\bar X + t_{1-\alpha/2,n-1}\frac{S}{\sqrt{n}}\right].
$$
где:
- $t_{n-1, 1-\frac{\alpha}{2}}$ — квантиль распределения Стьюдента с $n-1$ степенями свободы;
- $n$ — размер выборки;
- $\alpha$ — уровень значимости.
\\
\textbf{Доверительный интервал для $\sigma^2$ (уровень доверия $1-\alpha$)}

Для нормальной выборки:
$$
\frac{(n-1)S^2}{\sigma^2} \sim \chi^2_{n-1}.
$$
Отсюда доверительный интервал:
$$
\sigma^2 \in
\left[
\frac{(n-1)S^2}{\chi^2_{1-\alpha/2,n-1}},\;
\frac{(n-1)S^2}{\chi^2_{\alpha/2,n-1}}
\right].
$$
где:
- $(n-1)S^2 = \sum_{i=1}^n (X_i - \overline{X})^2$ — несмещённая оценка дисперсии;
- $\chi^2_{n-1, p}$ — квантиль распределения $\chi^2$ с $n-1$ степенями свободы.
\subsubsection{Поведение нормированных сумм для распределения Коши}
### Поведение нормированных сумм для распределения Коши

Пусть $X_i \sim C(a,b)$ — независимые одинаково распределённые случайные величины
(распределение Коши со сдвигом $a$ и масштабом $b$), а $S_n=\sum_{i=1}^n X_i$.

У распределения Коши характеристическая функция имеет вид
$$
\varphi(t)=\mathbb{E}e^{itX}=e^{ita-b|t|}.
$$

Тогда для нормированной суммы $\frac{S_n}{n}$:
$$
\varphi_{S_n/n}(t)=\mathbb{E}e^{itS_n/n}
= \left(\varphi\!\left(\frac{t}{n}\right)\right)^n
= \left(e^{i(t/n)a-b|t|/n}\right)^n
= e^{ita-b|t|}
= \varphi(t).
$$

Следовательно,
$$
\frac{S_n}{n} \sim C(a,b)\quad \forall n.
$$

Из-за тяжёлых хвостов распределения Коши отдельные наблюдения могут давать большой вклад в сумму,
поэтому сходимость $\frac{S_n}{n}$ к константе отсутствует (в частности, математическое ожидание не существует).

\subsection{Результаты:}
\subsubsection{ЗБЧ и ЦПТ}
\begin{center}
\includegraphics[width=0.7\linewidth]{мой отчёт/16.png}
\end{center}
\begin{center}
\includegraphics[width=0.7\linewidth]{мой отчёт/17.png}
\end{center}
Оба графика подтверждают нормальность распределения.

\subsubsection{Построение доверительных интервалов для матожидания и дисперсии}
\begin{center}
\includegraphics[width=0.7\linewidth]{мой отчёт/18.png}
\end{center}

\begin{center}
\includegraphics[width=0.7\linewidth]{мой отчёт/19.png}
\end{center}
\subsubsection{Поведение нормированных сумм}
\begin{center}
\includegraphics[width=0.7\linewidth]{мой отчёт/20.png}
\end{center}

На графике одной реализации видно, что величина $\frac{S_n}{n}$ при увеличении
n не стабилизируется и может резко изменяться. Такие скачки связаны с появлением отдельных больших значений, характерных для распределения Коши. Это наглядно показывает, что в данном случае среднее не сходится при росте числа слагаемых.

\section{Задание 6}
\subsection{Постановка задачи}
\begin{enumerate}
    
Вычислить следующий интеграл:

$$
\int_{-\infty}^\infty \int_{-\infty}^\infty \dots \int_{-\infty}^\infty 
\frac{e^{-\left(x_1^2 + \dots + x_{10}^2 + \frac{1}{27} \cdot x_1^2 \cdots x_{10}^2\right)}}{x_1^2 \cdots x_{10}^2} \, dx_1 dx_2 \dots dx_{10}
$$

\item Методом Монте-Карло

\item Методом квадратур, сводя задачу к вычислению собственного интеграла Римана.

\item Оценить точность вычислений для каждого из двух случаев.
\end{enumerate}

\subsection{Теоретические сведения}
\subsubsection{Метод Монте–Карло для вычисления интеграла}

Метод Монте–Карло основан на моделировании случайных величин и использовании выборочных средних для приближённого вычисления интересующих величин. В задачах численного интегрирования этот метод применяется следующим образом.

Пусть требуется вычислить интеграл
$$
\int_a^b f(x)\,\mathrm dx.
$$
Рассмотрим случайную величину $U$, равномерно распределённую на отрезке $[a,b]$. Тогда величина $f(U)$ также является случайной, и её математическое ожидание равно
$$
\mathbb E f(U)=\int_a^b f(x)\varphi(x)\,\mathrm dx,
$$
где $\varphi(x)=\frac{1}{b-a}$ — плотность распределения $U$.
Отсюда следует представление
$$
\int_a^b f(x)\,\mathrm dx=(b-a)\,\mathbb E f(U).
$$
Таким образом, значение интеграла можно приближённо найти, вычисляя выборочное среднее значений $f(U)$, полученных при моделировании равномерной случайной величины.Любой интеграл можно переписать как матожидание подходящей случайной величины.

Аналогичную идею применим к рассматриваемому многомерному интегралу. Пусть имеются случайные величины
$$
X_i\sim\mathcal{N(\mu=0,\sigma^2=\frac{1}{2})}\quad i=\overline{1,10},
$$
Тогда плотность каждой компоненты имеет вид
$$
p_i(x_i)=\frac{1}{\sqrt{\pi}}e^{-x_i^2},
$$
а совместная плотность вектора
$$
X=(X_1,\dots,X_{10})
$$
равна
$$
p(x)=\prod_{i=1}^{10}p_i(x_i)
=\frac{1}{\pi^5}\exp\!\left(-(x_1^2+\dots+x_{10}^2)\right),
\quad x\in\mathbb R^{10}.
$$

Введём обозначения
$$
y=\frac{1}{x_1^2\cdots x_{10}^2},
\qquad
f(x)=\pi^5\,y\,\exp\!\left(-\frac{y}{2^7}\right).
$$
Тогда исходный интеграл может быть представлен в виде математического ожидания:
$$
I=\int_{\mathbb R^{10}} f(x)\,p(x)\,\mathrm dx=\mathbb E f(X).
$$

Следовательно, для приближённого вычисления интеграла достаточно смоделировать $n$ независимых реализаций вектора $X$, вычислить соответствующие значения $f(X_k)$ и использовать выборочное среднее
$$
\overline X_n=\frac{1}{n}\sum_{k=1}^n f(X_k)
$$
в качестве оценки математического ожидания $\mathbb E f(X)$.

Для оценки погрешности воспользуемся неравенством Чебышёва:
$$
\mathbb P\bigl(|\overline X_n-\mathbb E X|\ge\varepsilon\bigr)
\le
\frac{\mathbb D X}{n\varepsilon^2}.
$$
Дисперсия $\mathbb D X$ аппроксимируется несмещённой выборочной дисперсией
$$
s_n^2=\frac{1}{n-1}\sum_{k=1}^n\bigl(f(X_k)-\overline X_n\bigr)^2.
$$
При достаточно большом размере выборки $(n\ge100)$ эта оценка обладает хорошей точностью. В результате получаем
$$
\mathbb P\bigl(|\overline X_n-\mathbb E X|\ge\varepsilon\bigr)
\le
\frac{s_n^2}{n\varepsilon^2}
=1-\alpha.
$$

В дальнейших вычислениях фиксируем $\varepsilon=1$ и $n=10^7$.


\subsubsection{Метод квадратур}


\textbf{Численное вычисление интеграла методом квадратур}

Для приближённого вычисления интеграла воспользуемся квадратурным методом, основанным на замене переменных и последующем применении метода прямоугольников.

Рассмотрим замену переменных
$$
x_i = \tan y_i,\qquad
\mathrm dx_i = \frac{\mathrm dy_i}{\cos^2 y_i},
\quad i=\overline{1,10}.
$$
При такой замене область интегрирования $\mathbb R^{10}$ переходит в гиперкуб
$[-\tfrac{\pi}{2},\,\tfrac{\pi}{2}]^{10}$.

После подстановки интеграл принимает вид
$$
\int_{-\frac{\pi}{2}}^{\frac{\pi}{2}}\!\!\dots
\int_{-\frac{\pi}{2}}^{\frac{\pi}{2}}
\exp\!\left(
-\sum_{i=1}^{10}\tan^2 y_i
-2^{-7}\prod_{i=1}^{10}\cot^2 y_i
\right)
\prod_{i=1}^{10}\frac{1}{\sin^2 y_i}\,
\mathrm dy_1\dots \mathrm dy_{10}.
$$

Обозначим
$$
y=(y_1,\dots,y_{10}),
$$
$$
f(y)=
\exp\!\left(
-\sum_{i=1}^{10}\tan^2 y_i
-2^{-7}\prod_{i=1}^{10}\cot^2 y_i
\right)
\prod_{i=1}^{10}\frac{1}{\sin^2 y_i}.
$$
Тогда интеграл можно записать в компактной форме:
$$
I=\int_{[-\frac{\pi}{2},\,\frac{\pi}{2}]^{10}} f(y)\,\mathrm dy.
$$

Поскольку подынтегральная функция является чётной по каждой переменной, область интегрирования может быть ограничена положительным ортантам:
$$
I = 2^{10}\int_{[0,\,\frac{\pi}{2}]^{10}} f(y)\,\mathrm dy.
$$

Для численного вычисления данного интеграла применим **метод средних прямоугольников**. Разобьём отрезок $[0,\tfrac{\pi}{2}]$ на $N$ равных частей длины
$$
h=\frac{\pi}{2N}.
$$
В качестве узлов квадратурной формулы используем середины отрезков:
$$
T=\left\{\frac{h}{2},\,\frac{3h}{2},\,\dots,\,\frac{(2N-1)h}{2}\right\}
=\{T_0,\dots,T_{N-1}\}.
$$

Переходя к декартовому произведению $T^{10}$, получаем равномерную 10-мерную сетку. Итоговая квадратурная формула имеет вид
$$
I \approx 2^{10}\,h^{10}\sum_{y\in T^{10}} f(y).
$$
Можно существенно сократить вычисления, заметив, что функция не меняет значение при перестановке аргументов. Тогда интеграл можно записать как
$$
    I = \sum_{0 \le k_1 \le k_2 \le \dots \le k_{10} \le N-1} \frac{N!}{n_1!\cdot \dots \cdot n_{N}!} f(T_{k_1}, \dots, T_{k_{10}}),
$$
где $n_i$ - число индексов, равных $i - 1, \quad i = \overline{1, N}$.
\\
\subsection{Оценка погрешности}

Рассмотрим оценку погрешности метода средних прямоугольников. Для этого разложим функцию
$f(\mathbf{x})$ в ряд Тейлора в окрестности центра $\mathbf{c}$ элементарного гиперкуба:
$$
f(\mathbf{x}) \approx
f(\mathbf{c})
+ \sum_{k=1}^n \frac{\partial f}{\partial x_k}\Big|_{\mathbf{c}}(x_k-c_k)
+ \frac{1}{2}\sum_{k,l=1}^n
\frac{\partial^2 f}{\partial x_k\partial x_l}\Big|_{\mathbf{c}}
(x_k-c_k)(x_l-c_l) + \dots
$$

В силу симметрии области интегрирования вклад линейных (нечётных) членов в интеграл равен нулю. Основной вклад в погрешность даёт квадратичная часть разложения. Интегрирование этих членов по одному гиперкубу приводит к оценке
$$
\text{погрешность на одном гиперкубе}
\;\sim\;
\frac{h^2}{24}
\sum_{k=1}^n
\frac{\partial^2 f}{\partial x_k^2}\Big|_{\mathbf{c}}
\cdot h^n.
$$

Пусть все вторые производные функции ограничены константой $M$. Тогда суммарная погрешность по всей области интегрирования оценивается сверху как
$$
E \sim
N^n \cdot \frac{h^2}{24} \cdot M \cdot h^n
= \frac{M}{24}\,(b-a)^n\,h^2.
$$

Несмотря на квадратичную зависимость от шага $h$, в многомерном случае возникает принципиальная проблема: число узлов сетки
$$
N^n = \left(\frac{b-a}{h}\right)^n
$$
растёт экспоненциально с увеличением размерности $n$. Это означает, что для достижения высокой точности требуется существенно уменьшать шаг разбиения, что приводит к резкому росту числа вычислений.

Таким образом, применение классических квадратурных методов в пространствах высокой размерности оказывается вычислительно неэффективным. Данное явление известно как *проклятие размерности*.

\subsection{Результаты:}
\subsubsection{Вычисление методом Монте-Карло}
Зафиксируем уровень значимости $N = 10^7$, $\epsilon = 1$ и проведём $n_{runs} = 15$ запусков

\begin{center}
\includegraphics[width=0.7\linewidth]{мой отчёт/21.png}
\end{center}
(1 - alpha) - вероятность вычисленного значения попасть в доверительный интервал\\ $[I - \varepsilon, I + \varepsilon].$.

\subsubsection{Вычисление методом квадратур}
Вычисление методом квадратур даёт значение $I_{quadrature} = 116.390258600636$.

\section{Задание 7}
\subsection{Постановка задачи}
\begin{enumerate}
    \item Методом случайного поиска найти минимальное значение функции $f$ на множестве $ A = \{x_1, x_2 : x_1^2 + x_2^2 \leq 1\} $, где
$$
f(x_1, x_2) = x_1^3 \sin \frac{1}{x_1} + 10x_1x_2^4 \cos \frac{1}{x_2}, \quad x_1, x_2 \neq 0
$$
При $x_1 = 0$ или $x_2 = 0$ функция доопределяется по непрерывности.

    \item Методом имитации отжига найти минимальное значение функции Розенброка $g$ в пространстве $\mathbb{R}^2$, где
$$
g(x) = (x_1 - 1)^2 + 100(x_2 - x_1^2)^2
$$

    \item Оценить точность и сравнить результаты со стандартными методами оптимизации.
\end{enumerate}
\subsection{Теоретические сведения}
\subsubsection{Метод случайного поиска}
Идея метода случайного поиска следующая: на каждом шаге из текущей точки делается случайная попытка перейти в новую точку на фиксированном расстоянии $r$.

Если новая точка:
\begin{enumerate}
    \item остаётся внутри допустимой области;
    \item даёт меньшее значение функции,
\end{enumerate}
то переход принимается. Иначе остаёмся в текущей точке.
В качестве критерия остановки будем использовать заранее заданное число итераций.

\subsubsection{Метод имитации отжига}
Название метода происходит от обжига в металлургии — это процесс, при котором металл постепенно охлаждатеся для изменения его физических свойств. Основные идеи метода состоят в следующем:\\
Новый кандидат формируется по формуле
$$
x_{k+1} = x_k + \mathcal{N}(0, \sigma^2 T_k).
$$

Если
$$
\Delta g = g(x_{k+1}) - g(x_k) \le 0,
$$
кандидат принимается.  
Иначе переход осуществляется с вероятностью
$$
p = \exp\left(-\frac{\Delta g}{T_k}\right).
$$

Температура понижается по закону
$$
T_k = \frac{T_0}{1 + \alpha k}.
$$
\subsection{Результаты:}
\subsubsection{Минимизация функции методом случайного поиска}
Применим метод случайного поиска для функции f. Зафиксируем $R = 1, niter = 100000$. Проведем по 5 вычислений для разных значений параметров.\\
при r = 0.1:
\begin{center}
\includegraphics[width=0.7\linewidth]{мой отчёт/22.png}
%r = результаты при r = 0.1
\end{center}
при r = 0.01:
\begin{center}
\includegraphics[width=0.7\linewidth]{мой отчёт/23.png}
%r = результаты при r = 0.01
\end{center}
при r = 0.001:
\begin{center}
\includegraphics[width=0.7\linewidth]{мой отчёт/24.png}
%r = результаты при r = 0.001
%
\end{center}

По результатам видно, что выбор шага заметно влияет на поведение алгоритма. При слишком маленьком шаге метод становится стабильнее, но чаще «застревает» в одном из локальных минимумов. При более крупном шаге он исследует область активнее, однако итоговые значения сильнее колеблются от запуска к запуску.

Для наглядности покажем путь случайного поиска:
\begin{center}
\includegraphics[width=0.7\linewidth]{мой отчёт/25.png}
\end{center}

\subsubsection{Минимизация функции методом имитации отжига}
Установим $T = 0.2, \alpha = 0.95, \sigma = 1, n = 100$. Количество запусков останется тем же.
\begin{center}
\includegraphics[width=0.7\linewidth]{мой отчёт/26.png}
\end{center}

\begin{center}
\includegraphics[width=0.7\linewidth]{мой отчёт/27.png}
\end{center}
Для наглядности покажем путь имитации отжига:
\begin{center}
\includegraphics[width=0.7\linewidth]{мой отчёт/28.png}
\end{center}
На графике можем увидеть, что метод действительно допускает ухудшения в отдельных итерациях при переходе к следующей точке.\\
Подводя итоги, метод имитации обжига работает намного лучше случайного поиска, однако даже он справляется с задачей не идеально. Возможно, при выборе других семейств функций алгоритм будет работать лучше.

\subsubsection{Сравнение со стандартными методами оптимизации}
Сравним работу случайных алгоритмов с классическим методом градиентного спуска. Градиент функции равен
$$ \nabla g(x_1, x_2) = \begin{bmatrix}
2 (x_1 - 1) - 400 (x_2 - x_1^2) x_1 \\
200 (x_2 - x_1^2)
\end{bmatrix} $$
Метод градиентного спуска чувствителен к начальному приближению, поэтому будем брать случайные начальные точки и запустим 1000 раз.
Получим следующие результаты:
Average function value: 7.735245388308443, variance: 258.568817998004
total best function value: 1.839364634030438e-10.\\
Градиентные методы демонстрируют более высокую скорость и точность на гладких функциях, однако стохастические методы обладают большей универсальностью и меньшей чувствительностью к начальному приближению.

\section{Задание 8}
\subsection{Постановка задачи}
\begin{enumerate}
    \item Применить метод Монте-Карло к решению первой краевой задачи для двумерного уравнения Лапласа в единичном круге:
$$
\begin{cases}
    \Delta u = 0, \quad (x,y) \in D = \{x, y: x^2+y^2 \leq 1\}, \\
    u|_{\partial D} = f(x, y), \\
    u\in C^2(D), \quad f\in C(\partial D), \\
\end{cases}
$$
    \item Для функции $f(x, y) = x^2-y^2$ найти аналитическое решение и сравнить с полученным по методу Монте-Карло.
\end{enumerate}
\subsection{Теоретические сведения}

\textbf{Идея вероятностного подхода}

Для гармонической функции значение в  mвнутренней точке области связано со случайным блужданием: если многократно запускать траекторию из выбранного узла и останавливать её в момент первого попадания на границу, то среднее значение граничной функции в точках выхода будет приближать искомое решение.

В дискретной постановке это означает следующее:
- внутри круга строится сетка;
- из каждого внутреннего узла запускается серия случайных блужданий;
- на каждом шаге траектория переходит в один из соседних узлов с одинаковой вероятностью;
- после достижения границы фиксируется значение $f$;
- итоговое значение $u$ в стартовом узле берётся как среднее по всем траекториям.
\\
\textbf{Сеточная аппроксимация области}

Область расчёта задаётся на равномерной прямоугольной сетке в квадрате $[-1,1]\times[-1,1]$. После этого из всех узлов оставляются только те, для которых выполняется условие
$$
x^2+y^2\le 1.
$$

Среди них выделяются:
- внутренние узлы, из которых можно сделать шаг во все четыре стороны, не покидая область;
- граничные узлы, для которых хотя бы один сосед выходит за пределы круга.

На граничных узлах решение задаётся сразу:
$$
u(x,y)=f(x,y).
$$

\textbf{Схема случайного блуждания}

Используется простейшее блуждание на декартовой сетке. Из текущего узла частица переходит в один из четырёх соседних узлов:
- вверх,
- вниз,
- влево,
- вправо,

причём вероятность каждого перехода равна $\frac14$.

Одно моделирование устроено так:
1. выбирается стартовый внутренний узел;
2. случайно определяется направление очередного шага;
3. траектория переносится в соседний узел;
4. если достигнута граница, процесс останавливается;
5. если граница ещё не достигнута, шаги продолжаются;
6. в момент остановки записывается значение граничной функции в точке выхода.
\\
\textbf{Формула Монте-Карло}

Пусть из точки $(x,y)$ выполнено $N$ независимых блужданий, а $(x_k^{*},y_k^{*})$ — точка первого выхода $k$-й траектории на границу. Тогда приближение решения задаётся формулой
$$
u_N(x,y)=\frac{1}{N}\sum_{k=1}^{N} f(x_k^{*},y_k^{*}).
$$

При увеличении числа траекторий значение $u_N(x,y)$ стабилизируется, а статистическая погрешность убывает как величина порядка $N^{-1/2}$.\\

\subsection{Результаты:}
\subsubsection{Решение методом Монте-Карло}
Установим $N_iter$ - количество пробегов для каждой внутренней точки и
$N$ - количество узлов по каждой оси \((x, y)\)
\begin{center}
\includegraphics[width=0.7\linewidth]{мой отчёт/29.png}
\end{center}
\subsubsection{Аналитическое решение для $f(x,y)=x^2-y^2$}

Рассмотрим функцию
$$
u_{\text{exact}}(x,y)=x^2-y^2.
$$

Проверим, что она гармоническая в круге:
$$
\Delta u_{\text{exact}}=\frac{\partial^2}{\partial x^2}(x^2-y^2)+\frac{\partial^2}{\partial y^2}(x^2-y^2)=2-2=0.
$$

На границе области функция принимает значения
$$
u_{\text{exact}}|_{\partial D}=x^2-y^2=f(x,y),
$$
то есть точно удовлетворяет заданному граничному условию. Следовательно, искомое аналитическое решение задачи имеет вид
$$
u(x,y)=x^2-y^2.
$$
\begin{center}
\includegraphics[width=0.7\linewidth]{мой отчёт/30.png}
\end{center}
\subsubsection{Сравнение численного и точного решения}
\begin{center}
\includegraphics[width=0.7\linewidth]{мой отчёт/31.png}
\end{center}

\section{Задание 9}
\subsection{Постановка задачи}
Рассмотреть два вида гауссовских процессов:

- Винеровский процесс $ W(t), t \in [0; 1], W(0) = 0 $

- Процесс Орнштейна–Уленбека $ X(t), t \in [0; 1], X(0) = X_0 $, т.е. стационарный марковский гауссовский процесс. Начальные значения $ X_0 $ следует выбирать случайным образом так, чтобы полученный процесс был стационарным.

Для данных процессов:
\begin{enumerate}
    \item Найти ковариационную функцию и переходные вероятности.
    \item Промоделировать независимые траектории процесса с данными переходными вероятностями методом добавления разбиения отрезка.
    \item Построить график траектории, не соединяя точки ломаной, с целью получения визуально непрерывной линии.
\end{enumerate}
\subsection{Теоретические сведения}
\subsubsection{Винеровский процесс}

Винеровским процессом $W(t),\, t \geqslant 0$, называется случайный процесс, для которого верно:
1. $W(0) = 0 \text{ п.н.}$
2. $W(t)$ &mdash; $\text{ процесс с независимыми приращениями, то есть } \forall t_0, t_1, \dots, t_n \geqslant 0 \text{ таких, что } 0 = t_0 < t_1 < \dots < t_n, \text{ случайные величины } X_{t_0}, X_{t_1} - X_{t_0},\dots, X_{t_n} - X_{t_{n-1}} \text{ независимы}.$
3. $W(t) - W(s) \sim N(0, \sigma^2(t-s)), \forall 0 \leqslant s < t,$ где $N(0, \sigma^2(t-s))$ &mdash; нормальное распределение со средним 0 и дисперсией $\sigma^2(t-s)$.

Винеровский процесс является марковским и гауссовским.
\\
Найдем ковариационную функцию: Пусть $s \leqslant t$, тогда
$$\text{cov}(W(s), W(t)) = \mathbb{E}(W(s) - \mathbb{E}W(s))(W(t) - \mathbb{E}W(t)) = \mathbb{E}W(s)W(t) = \mathbb{E}W(s)((W(t) - W(s)) + W(s)) = \mathbb{E}W(s)(W(t) - W(s)) + \mathbb{E} W^2(s).$$
В силу независимости приращений получаем, что случайные величины
$$W(s) = W(s) - W(0) \text{ и } W(t) - W(s)$$
независимы, поэтому
$$\mathbb{E}W(s)(W(t) - W(s)) = \mathbb{E}W(s) \cdot \mathbb{E}(W(t) - W(s)) = 0.$$
Таким образом,
$$\text{cov}(W(s), W(t)) = \mathbb{E}(W^2(s)) = \text{Var}(W^2(s)) = s\sigma^2.$$
В общем случае получаем
$$\text{cov}(W(s), W(t)) = \text{min}(s, t)\sigma^2.$$
Теперь рассмотрим переходные вероятности. Пусть, как и раньше, $0 \leqslant s < t$. Нас интересует условная плотность распределения:
$$p_{W(t)|W(s)}(y|x) = \frac{p_{W(s),W(t)}(x,y)}{p_{W(s)}(x)}$$
Заметим, что
$$W(t) = W(t) - W(0) \sim N(0, \sigma^2t),\, W(s) = W(s) - W(0) \sim N(0, \sigma^2s),$$
поэтому случайный вектор $(W(s), W(t))$ имеет двумерное нормальное распределение с матожиданием $\mu = (0, 0)$ и ковариационной матрицей
$$\Sigma = \begin{bmatrix}
\text{cov}(W(s), W(s)) & \text{cov}(W(s), W(t))\\
\text{cov}(W(t), W(s)) & \text{cov}(W(t), W(t))\\
\end{bmatrix} = \sigma^2\begin{bmatrix}
s & s\\
s & t\\
\end{bmatrix},$$
$$\text{det}(\Sigma) = \sigma^4s(t-s),$$
$$\Sigma^{-1} = \frac{1}{\sigma^2s(t-s)}\begin{bmatrix}
t & -s\\
-s & s\\
\end{bmatrix}.$$
Таким образом,
$$
p_{W(s), W(t)}(x, y) = \frac{1}{2\pi \sigma^2 \sqrt{s(t-s)}}e^{-\frac{1}{2\sigma^2s(t-s)}\begin{bmatrix}x & y\end{bmatrix}\begin{bmatrix}
t & -s\\
-s & s\\
\end{bmatrix}\begin{bmatrix}x \\ y\end{bmatrix}} =
\frac{1}{2\pi \sigma^2 \sqrt{s(t-s)}}e^{-\frac{1}{2\sigma^2s(t-s)}(tx^2 - 2sxy + sy^2)},
$$
$$p_{W(s)}(x) = \frac{1}{\sqrt{2\pi s}\sigma}e^{-\frac{1}{2}\frac{x^2}{\sigma^2s}},$$
и
$$
p_{W(t)|W(s)}(y|x) = \frac{\sqrt{2\pi s}\sigma}{2\pi \sigma^2 \sqrt{s(t-s)}}e^{-\frac{1}{2\sigma^2s(t-s)}(tx^2 - 2sxy + sy^2) + \frac{1}{2}\frac{x^2}{\sigma^2s}} =
\frac{1}{\sqrt{2\pi(t-s)} \sigma}e^{-\frac{1}{2\sigma^2s(t-s)}(sx^2 - 2sxy + sy^2)} =
\frac{1}{\sqrt{2\pi(t-s)} \sigma}e^{-\frac{1}{2\sigma^2(t-s)}(y - x)^2}.
$$
Значит,
$$(W(t)|W(s)=x) \sim N(x, \sigma^2(t - s))$$

Для метода разбиения отрезка нам также понадобится условная вероятность для средней точки отрезка при известных крайних. Пусть $0 \leqslant t_1 < t_2 < t_3$. Тогда
$$p_{W(t_2)|W(t_1),W(t_3)}(y|x,z) = \frac{p_{W(t_1),W(t_2),W(t_3)}(x,y,z)}{p_{W(t_1),W(t_3)}(x,z)},$$
$$\Sigma_{t_1, t_2, t_3} = \sigma^2\begin{bmatrix}
t_1 & t_1 & t_1\\
t_1 & t_2 & t_2\\
t_1 & t_2 & t_3\\
\end{bmatrix},$$
$$\text{det}(\Sigma_{t_1, t_2, t_3}) = \sigma^6t_1(t_2 - t_1)(t_3 - t_2),$$
$$\Sigma_{t_1, t_2, t_3}^{-1} = \frac{1}{\sigma^2t_1(t_2 - t_1)(t_3 - t_2)}\begin{bmatrix}
t_2(t_3 - t_2) & t_1(t_2 - t_3) & 0\\
t_1(t_2 - t_3) & t_1(t_3 - t_1) & t_1(t_1 - t_2)\\
0 & t_1(t_1 - t_2) & t_1(t_2 - t_1)\\
\end{bmatrix},$$
$$p_{W(t_1),W(t_2),W(t_3)}(x,y,z) = \frac{1}{\left(2\pi\sigma^2\right)^{\frac{3}{2}}\sqrt{t_1(t_2 - t_1)(t_3 - t_2)}}e^{-\frac{1}{2\sigma^2t_1(t_2 - t_1)(t_3 - t_2)}\begin{bmatrix}x & y & z\end{bmatrix}\begin{bmatrix}
t_2(t_3 - t_2) & t_1(t_2 - t_3) & 0\\
t_1(t_2 - t_3) & t_1(t_3 - t_1) & t_1(t_1 - t_2)\\
0 & t_1(t_1 - t_2) & t_1(t_2 - t_1)\\
\end{bmatrix}\begin{bmatrix}x \\ y \\ z\end{bmatrix}}.$$

В методе разбиения отрезка мы будем брать $t_2 = \frac{t_1 + t_3}{2}$, поэтому
$$\text{det}(\Sigma_{t_1, t_2, t_3}) = \sigma^6t_1(\frac{t_1 + t_3}{2} - t_1)(t_3 - \frac{t_1 + t_3}{2}) = \sigma^6t_1\left(\frac{t_3 - t_1}{2}\right)^2,$$
$$\Sigma_{t_1, t_2, t_3}^{-1} = \frac{1}{\sigma^2t_1\left(\frac{t_3 - t_1}{2}\right)^2}\begin{bmatrix}
\frac{t_1 + t_3}{2}\frac{t_3 - t_1}{2} & t_1\frac{t_1 - t_3}{2} & 0\\
t_1\frac{t_1 - t_3}{2} & t_1(t_3 - t_1) & t_1\frac{t_1 - t_3}{2}\\
0 & t_1\frac{t_1 - t_3}{2} & t_1\frac{t_3 - t_1}{2}\\
\end{bmatrix},$$
$$p_{W(t_1),W(t_2),W(t_3)}(x,y,z) = \frac{1}{\left(2\pi\sigma^2\right)^{\frac{3}{2}}\sqrt{t_1}\frac{t_3 - t_1}{2}}e^{-\frac{2}{\sigma^2t_1(t_3 - t_1)^2}\begin{bmatrix}x & y & z\end{bmatrix}\begin{bmatrix}
\frac{t_3^2 - t_1^2}{4} & t_1\frac{t_1 - t_3}{2} & 0\\
t_1\frac{t_1 - t_3}{2} & t_1(t_3 - t_1) & t_1\frac{t_1 - t_3}{2}\\
0 & t_1\frac{t_1 - t_3}{2} & t_1\frac{t_3 - t_1}{2}\\
\end{bmatrix}\begin{bmatrix}x \\ y \\ z\end{bmatrix}} =
\frac{1}{\pi^{\frac{3}{2}}\sigma^3\sqrt{2 t_1}(t_3 - t_1)}e^{-\frac{2}{\sigma^2t_1(t_3 - t_1)^2}\left(\frac{t_3^2 - t_1^2}{4}x^2 + y(x + z - y)t_1(t_1 - t_3) + t_1\frac{t_3 - t_1}{2}z^2\right)},
$$
$$
p_{W(t_1), W(t_3)}(x, z) = \frac{1}{2\pi \sigma^2 \sqrt{t_1(t_3-t_1)}}e^{-\frac{1}{2\sigma^2t_1(t_3-t_1)}(t_3x^2 - 2t_1xz + t_1z^2)},
$$
$$p_{W(t_2)|W(t_1),W(t_3)}(y|x,z) = \frac{2\pi \sigma^2 \sqrt{t_1(t_3-t_1)}}{\pi^{\frac{3}{2}}\sigma^3\sqrt{2 t_1}(t_3 - t_1)}e^{-\frac{2}{\sigma^2t_1(t_3 - t_1)^2}\left(\frac{t_3^2 - t_1^2}{4}x^2 + y(x + z - y)t_1(t_1 - t_3) + t_1\frac{t_3 - t_1}{2}z^2\right) + \frac{1}{2\sigma^2t_1(t_3-t_1)}(t_3x^2 - 2t_1xz + t_1z^2)}.$$
$$ -\frac{2}{\sigma^2t_1(t_3 - t_1)^2}\left(\frac{t_3^2 - t_1^2}{4}x^2 + y(x + z - y)t_1(t_1 - t_3) + t_1\frac{t_3 - t_1}{2}z^2\right) + \frac{1}{2\sigma^2t_1(t_3-t_1)}(t_3x^2 - 2t_1xz + t_1z^2) = $$
$$ = \frac{1}{2\sigma^2t_1(t_3 - t_1)^2}\left(-(t_3 - t_1)(t_3 + t_1)x^2 + 4t_1(t_3 - t_1)y(x + z - y) - 2t_1(t_3 - t_1)z^2 + t_3(t_3 - t_1)x^2 - 2t_1(t_3 - t_1)xz + t_1(t_3 - t_1)z^2\right) = $$
$$ = \frac{1}{2\sigma^2t_1(t_3 - t_1)^2}\left(-4t_1(t_3 - t_1)y^2 + 4t_1(t_3 - t_1)(x + z)y + (-(t_3 - t_1)(t_3 + t_1)x^2 - 2t_1(t_3 - t_1)z^2 + t_3(t_3 - t_1)x^2 - 2t_1(t_3 - t_1)xz + t_1(t_3 - t_1)z^2)\right) = $$
$$ = -\frac{1}{2\sigma^2(t_3 - t_1)}\left(4y^2 - 4y(x + z) + (x^2 + 2xz + z^2)\right) = -\frac{2}{\sigma^2(t_3 - t_1)}\left(y - \frac{x + z}{2}\right)^2 \Rightarrow$$
$$ \Rightarrow p_{W(t_2)|W(t_1),W(t_3)}(y|x,z) = \frac{1}{\sigma\sqrt{\pi \frac{t_3-t_1}{2}}}e^{-\frac{2}{\sigma^2(t_3 - t_1)}\left(y - \frac{x + z}{2}\right)^2} = \frac{1}{\sigma\sqrt{2\pi \frac{t_3-t_1}{4}}}e^{-\frac{1}{2\sigma^2\frac{t_3 - t_1}{4}}\left(y - \frac{x + z}{2}\right)^2}.$$
Таким образом, получаем, что
$$(W(t_2)|W(t_1)=x,W(t_3)=z) \sim N(\frac{x + z}{2}, \sigma^2\frac{t_3 - t_1}{4}).$$
\subsubsection{Процесс Орнштейна-Уленбека}

Процесс Орнштейна-Уленбека $X(t), t \in [0, 1], X(0) = X_0$ &mdash; единственный нетривиальный гауссовский процесс, являющийся стационарным и обладающий марковским свойством. Начальные значения $X_0$ выбираются случайным образом так, чтобы процесс был стационарным. Для этого необходимо и достаточно, чтобы $X_0 \sim N(\mu, \sigma^2)$
\newline

Выведем ковариационную функцию процесса. Пусть $0 \leqslant s < t$. Без ограничения общности, рассмотрим процессы с $\mu = 0$. В силу стационарности процесса,
$$\text{cov}(X(s), X(t)) = \text{cov}(X(0), X(t-s)) = K(t-s).$$
Рассмотрим условное матожидание $\mathbb{E}(X(t + h)|X(t))$. В силу того, что $X(t + h)$ и $X(t)$ mdash; совместно гауссовские случайные величины,
$$\mathbb{E}(X(t + h)|X(t)) = \frac{\text{cov}(X(t + h), X(t))}{\text{Var(X(t))}}X(t) = \frac{K(h)}{\sigma^2}X(t).$$
Теперь посчитаем величину $\mathbb{E}(X(t + h + s)|X(t))$ двумя способами. С одной стороны,
$$\mathbb{E}(X(t + h + s)|X(t)) = \frac{K(h + s)}{\sigma^2}X(t).$$
С другой же, используя марковское свойство и свойства условного матожидания,
$$\mathbb{E}(X(t + h + s)|X(t)) = \mathbb{E}(\mathbb{E}(X(t + h + s)|X(t + h))|X(t)) = \mathbb{E}(\frac{K(s)}{\sigma^2}X(t + h)|X(t)) = \frac{K(s)}{\sigma^2}\frac{K(h)}{\sigma^2}X(t).$$
Таким образом,
$$\frac{K(h + s)}{\sigma^2}X(t) = \frac{K(s)}{\sigma^2}\frac{K(h)}{\sigma^2}X(t),$$
$$\frac{K(h + s)}{\sigma^2} = \frac{K(s)}{\sigma^2}\frac{K(h)}{\sigma^2}.$$
В силу непрерывности решений, получаем, что уравнение имеет единственное решение
$$\frac{K(h)}{\sigma^2} = e^{-\theta h}$$
для некоторой константы $\theta$. В силу того, что корреляция должна быть ограничена, получаем, что $\theta > 0$. В общем случае, имеем
$$K(h) = \sigma^2 e^{\theta\vert h \vert}$$
Процесс нахождения переходных вероятностей аналогичен случаю винеровского процесса, в котором нужно заменить ковариационную функцию на соответствующую функцию для процесса Орнштейна-Уленбека, и учесть ненулевое матожидание. В итоге получим, что если $X_0 \sim N(\mu, \sigma^2)$, то
$$(X(t) | X_0 = x) \sim N(\mu + (x - \mu)e^{-\theta t}, \sigma^2(1 - e^{-2 \theta t})),$$
$$(X(t_2) | X(t_1) = x,X(t_3) = z) \sim N(\frac{x + z}{e^{\theta h} + e^{-\theta h}}, \sigma^2\frac{e^{2 \theta h} - 1}{e^{2 \theta h} + 1}),$$
где $t \geqslant 0,\, t_1 < t_3,\, t_2 = \frac{t_1 + t_3}{2},\, h = \frac{t_3 - t_1}{2}$.
\subsubsection{Метод разбиения отрезка}
Для моделирования независимых траекторий гауссовского марковского процесса удобно использовать \textbf{метод добавления разбиения отрезка}. Его идея состоит в том, что сначала задаются значения процесса в небольшом числе узлов, после чего разбиение временного интервала последовательно уточняется: между уже построенными точками добавляются новые моменты времени, а значения процесса в этих точках генерируются по **условным переходным распределениям**. Такой подход позволяет получать траектории, согласованные с конечномерными распределениями процесса.

Пусть временной интервал равен $[0,1]$. На первом шаге выбирается грубое разбиение, например только начальная и конечная точки. Затем на каждом этапе каждый промежуток $[t_i,t_{i+1}]$ делится пополам, и в середине $t_{mid}=\frac{t_i+t_{i+1}}{2}$ моделируется значение процесса с использованием условного распределения при известных значениях в концах промежутка. После этого процедура повторяется для новых отрезков. В результате после нескольких шагов уточнения получается достаточно плотная сетка, на которой строится траектория процесса.

Будем реализовывать метод через переходные вероятности процесса, то есть каждое новое значение генерируется не произвольно, а в соответствии с распределением
$$
X(t_{k+1}) \mid X(t_k).
$$
Так как рассматриваемые процессы являются гауссовскими, это условное распределение является нормальным, а значит для моделирования достаточно знать его математическое ожидание и дисперсию.
\\
\textbf{Винеровский процес}

Для винеровского процесса \(W(t)\) выполняется
$$
W(t_{k+1}) - W(t_k) \sim \mathcal{N}(0,\, t_{k+1}-t_k).
$$
Поэтому при добавлении новой точки траектория достраивается по формуле
$$
W(t_{k+1}) = W(t_k) + \xi_k, \qquad \xi_k \sim \mathcal{N}(0,\, \Delta t_k),
$$
где $\Delta t_k = t_{k+1}-t_k$.

Таким образом, на каждом новом участке добавляется независимое нормальное приращение с дисперсией, равной длине соответствующего временного промежутка. Это и позволяет последовательно уточнять разбиение и получать реализацию винеровского процесса.
\\
\textbf{Процесс Орнштейна-Уленбека}

Для процесса Орнштейна-Уленбека переходное распределение имеет вид
$$
X(t_{k+1}) \mid X(t_k) \sim \mathcal{N}\left(X(t_k)e^{-\theta \Delta t_k},\,
\frac{\sigma^2}{2\theta}\left(1-e^{-2\theta \Delta t_k}\right)\right).
$$
Следовательно, при добавлении новой точки используется формула
$$
X(t_{k+1}) = X(t_k)e^{-\theta \Delta t_k} + \eta_k,
$$
где
$$
\eta_k \sim \mathcal{N}\left(0,\,
\frac{\sigma^2}{2\theta}\left(1-e^{-2\theta \Delta t_k}\right)\right).
$$

Начальное значение $X(0)$ выбирается из стационарного распределения
$$
X(0)\sim \mathcal{N}\left(0,\frac{\sigma^2}{2\theta}\right),
$$
что обеспечивает стационарность всей построенной траектории.
\\
\textbf{Итоговый алгорим выглядит так:}
\begin{enumerate}
    \item Задаётся начальное разбиение отрезка времени.
    \item В начальных узлах моделируются значения процесса.
    \item Каждый промежуток делится на более мелкие части.
    \item В новых узлах значения процесса генерируются по соответствующим переходным вероятностям.
    \item Процедура повторяется до получения достаточно плотной сетки.
\end{enumerate}

Преимущество метода добавления разбиения состоит в том, что он позволяет строить траекторию постепенно, сохраняя вероятностную структуру процесса на каждом этапе уточнения. Поэтому полученные реализации корректно отражают свойства винеровского процесса и процесса Орнштейна-Уленбека.

\subsection{Результаты:}
\subsubsection{Моделирование процессов}

\begin{center}
\includegraphics[width=0.7\linewidth]{мой отчёт/32.png}
\end{center}
На гррафике можем заметить особенность винеровского процесса - тренд на увеличение разброс при увеличени рассматриваемого промежутка времени.\\
при $\theta=0.7, \mu=1.0, \sigma=0.3, depth=10$:
\begin{center}
\includegraphics[width=0.7\linewidth]{мой отчёт/33.png}
\end{center}
при $theta=100, mu=0, sigma=1, depth=15$:
\begin{center}
\includegraphics[width=0.7\linewidth]{мой отчёт/33.png}
\end{center}
Аналогично, можем заметить особенность процесса Орнштейна-Уленбека - траектории болтаются вокруг уровня $\mu$ подобно пружине, что демонстрирует наличие памяти.
\section{Задание 10}
\subsection{Постановка задачи}
Произвести фильтрацию одномерного процесса Орнштейна–Уленбека:

1. Используя генератор белого шума, добавить к реализации процесса Орнштейна–Уленбека случайную ошибку с заранее известной дисперсией.
2. При помощи одномерного фильтра Калмана оценить траекторию процесса по зашумленному сигналу, считая известными параметры шума и процесса.
3. Рассмотреть следующие виды шума:
\begin{enumerate}
    \item Гауссов
    \item Коши (шум имеет распределение Коши)
\end{enumerate}
\subsection{Теоретические сведения}
\textbf{Дискретная модель процесса}

Процесс Орнштейна–Уленбека в дискретном виде представляется линейной динамической системой:

$$
x_{n+1} = a x_n + v_n, \quad v_n \sim \mathcal{N}(0, q),
$$

где:
- $a$ — коэффициент динамики,
- $v_n$ — белый шум процесса,
- $q$ — дисперсия шума процесса.

Начальное значение:
$$
x_1 \sim \mathcal{N}(0, \sigma^2).
$$

Наблюдаемый процесс задаётся как:
$$
y_n = x_n + \varepsilon_n, \quad \varepsilon_n \sim \mathcal{N}(0, r),
$$

где $r$ — дисперсия шума измерений (считается известной)
\newpage

\textbf{Связь параметров с процессом Орнштейна–Уленбека}

Используя ковариационную функцию процесса, получаем:

$$
\mathrm{Var}(x_n) = \sigma^2,
$$

$$
\mathrm{Cov}(x_n, x_{n+1}) = \sigma^2 e^{-\lambda h},
$$

$$
\mathrm{Var}(x_{n+1}) = \sigma^2.
$$

С другой стороны, из динамической модели:

$$
\mathrm{Cov}(x_n, x_{n+1}) = a \sigma^2,
$$

$$
\mathrm{Var}(x_{n+1}) = a^2 \sigma^2 + q.
$$

Отсюда получаем параметры модели:

$$
a = e^{-\lambda h}, \quad q = \sigma^2 (1 - e^{-2\lambda h}).
$$

\subsubsection{Фильтр Калмана}

Для оценки скрытого состояния используется одномерный фильтр Калмана.

Обозначим:
- $\hat{x}_{n|n}$ — оценка состояния на шаге $n$,
- $\hat{x}_{n|n-1}$ — прогноз,
- $P_{n|n}$ — дисперсия ошибки оценки,
- $P_{n|n-1}$ — прогнозируемая дисперсия.

\textbf{Шаг 1.} Прогнозируем значение процесса и дисперсию ошибки

$$
\hat{x}_{n|n-1} = a \hat{x}_{n-1|n-1},
$$

$$
P_{n|n-1} = a^2 P_{n-1|n-1} + q.
$$

\textbf{Шаг 2.} Вычисляем разницу между наблюдаемым процессом и прогнозом и коэффициент усиления

$$
\delta_n = y_n - \hat{x}_{n|n-1}.
$$

$$
K_n = \frac{P_{n|n-1}}{P_{n|n-1} + r}.
$$

\textbf{Шаг 3.}В качестве результата фильтрации берем выпуклую комбинацию наблюдаемого и предсказанного значения:
$$
\hat{x}_{n|n} = \hat{x}_{n|n-1} + K_n \delta_n,
$$

$$
P_{n|n} = (1 - K_n) P_{n|n-1}.
$$

\newpage
\textbf{Доверительные интервалы}

При гауссовских предположениях ошибка фильтрации имеет нормальное распределение.

Доверительный интервал уровня $\alpha$ задаётся как:

$$
\hat{x}_{n|n} \pm \sqrt{P_{n|n}} \cdot z_{\alpha/2},
$$

где $z_{\alpha/2}$ — квантиль стандартного нормального распределения.
\newline

\textbf{Особенности негауссовского шума}

Если шум имеет распределение Коши, то:

- дисперсия не существует,
- предположения фильтра Калмана нарушаются.

В этом случае фильтр Калмана перестаёт быть оптимальным, и качество оценки может существенно ухудшаться.

\subsection{Результаты:}
\subsubsection{Фильтр Калмана для процесса О-У с Гауссовским шумом}
\begin{center}
\includegraphics[width=0.7\linewidth]{мой отчёт/36.png}
\end{center}
\begin{center}
\includegraphics[width=0.7\linewidth]{мой отчёт/37.png}
\end{center}
\begin{center}
\includegraphics[width=0.7\linewidth]{мой отчёт/38.png}
\end{center}
Таким образом, фильтр Калмана позволяет эффективно восстанавливать траекторию процесса Орнштейна–Уленбека при гауссовском шуме.

\subsubsection{Фильтр Калмана для процесса О-У с шумом Коши}
\begin{center}
\includegraphics[width=0.7\linewidth]{мой отчёт/39.png}
\end{center}

\begin{center}
\includegraphics[width=0.7\linewidth]{мой отчёт/40.png}
\end{center}
На графике невязки для шума Коши наблюдаются редкие, но очень большие пики. Это связано с тяжёлыми хвостами распределения Коши и наличием выбросов в наблюдениях. В отличие от гауссовского шума, для распределения Коши не существует дисперсии, поэтому предпосылки применения фильтра Калмана нарушаются. Следовательно, при отклонении от гауссовских предположений фильтр Калмана может давать нестабильные результаты и требует осторожной интерпретации.

\section{Задание 11}
\subsection{Постановка задачи}
Построить двумерное пуассоновское поле, отвечающее сложному пуассоновскому процессу:
\begin{enumerate}
    \item Система массового обслуживания.
   Первая координата поля — время поступления заявки в СМО (распределенное равномерно), а вторая — время обслуживания заявки (распределение $ \chi^2 $ с десятью степенями свободы).
   \item Система массового обслуживания с циклической интенсивностью $ \lambda(1 + \cos(t)) $ и единичными скачками. 
   При помощи метода Льюиса и Шедлеара, свести задачу моделирования неоднородного пуассоновского процесса к моделированию двумерного пуассоновского поля, где первая координата распределена равномерно, а вторая имеет распределение Бернулли.
   \item Работа страховой компании: 
   Первая координата — момент наступления страхового случая (равномерное распределение), вторая — величина ущерба (распределение Парето). Поступление капитала считать линейным по времени со скоростью $ c > 0 $, начальный капитал $ W > 0 $.
\end{enumerate}

\subsection{Теоретические сведения}
\subsubsection{Система массового обслуживания}

Пусть на интервале $[0, T]$ поступает случайное число заявок:
$$
N \sim \text{Pois}(\lambda T).
$$

Моменты поступления заявок $t_1,\dots,t_N$ будем считать независимыми и равномерно распределёнными на отрезке [0,T]:
$$
t_i \sim U[0, T].
$$

Время обслуживания каждой заявки моделируется случайными величинами:
$$
\tau_i \sim \chi^2(10),
$$
которые также считаются независимыми.

Таким образом, каждая заявка задаётся точкой $(t_i, \tau_i)$, что формирует двумерное пуассоновское поле.

\subsubsection{Система с переменной интенсивностью}

Интенсивность:
$$
\lambda(t) = \lambda_0 (1 + \cos t).
$$

Так как процесс неоднородный, используется метод Льюиса–Шедлера.

Идея метода:

1. Выбирается максимальная интенсивность:
$$
\lambda_{\max} = 2\lambda_0.
$$

2. Генерируется однородный пуассоновский процесс.

3. Каждое событие принимается с вероятностью:
$$
p(t) = \frac{\lambda(t)}{\lambda_{\max}} = \frac{1 + \cos t}{2}.
$$

Отброшенные события удаляются.

Вторая координата моделируется распределением Бернулли:
$$
X_i \sim \text{Bernoulli}(p).
$$
\subsubsection{Модель страховой компании}

Набор точек $(t_i, \xi_i)$ образует двумерное пуассоновское поле,
где первая координата — момент страхового случая, вторая — величина ущерба.$\newline$
Будем моделировать динамику капитала компании, учитывая случайные страховые случаи с убытками, распределёнными по закону Парето, и линейное поступление капитала.


1. \textbf{Убытки}: страховые случаи генерируются как пуассоновский процесс с интенсивностью $\lambda$. Размер убытков ($\tau$) следует распределению Парето:
   $$
   \tau = \frac{b}{U^{1/a}}, \quad U \sim \text{Uniform}(0, 1),
   $$
   где $a$ — параметр формы, $b$ — параметр масштаба.
2. \textbf{Поступления}: капитал увеличивается линейно со скоростью $c$.
3. \textbf{Капитал}: рассчитывается как:
   $$
   W(t) = W_0 + ct - \sum_{i=1}^{N(t)} \xi_i,
   $$
   где $W_0$ — начальный капитал. При $W(t) < 0$ капитал обнуляется (банкротство).

Модель позволяет анализировать устойчивость компании при разных параметрах и оценивать риск банкротства.

где $N(t)$ — пуассоновский процесс, $\xi_i$ — независимые выплаты.

\subsection{Результаты:}
При $\lambda = 1$ и $T = 50$:
\subsubsection{Пуассоновское поле для СМО}
\begin{center}
\includegraphics[width=0.7\linewidth]{мой отчёт/41.png}
\end{center}

\subsubsection{Пуассоновское поле для системы с переменной интенсивностью}
При $l_0 = 75$ и $T = 50$:
\begin{center}
\includegraphics[width=0.7\linewidth]{мой отчёт/42.png}
\end{center}

Для анализа полученного процесса была построена гистограмма распределения моментов событий по времени.

\begin{center}
\includegraphics[width=0.7\linewidth]{мой отчёт/43.png}
\end{center}

Гистограмма позволяет оценить, как изменяется плотность событий на интервале времени. В отличие от scatter-графика, она наглядно показывает, в какие моменты времени события происходят чаще, а в какие реже.

Из полученного графика видно, что число событий распределено неравномерно: наблюдаются выраженные области с повышенной и пониженной концентрацией точек.

Такая структура соответствует заданной функции интенсивности:
$λ(t) = λ₀ (1 + \cos(t))$

Максимумы на гистограмме приходятся на значения $t$, где $\cos(t)$ близок к 1, а минимумы — где $\cos(t)$ близок к -1.

Таким образом, построенная гистограмма подтверждает корректность моделирования неоднородного пуассоновского процесса методом прореживания.
\subsubsection{Модель страховой компании}
Зафиксируем различные входные параметры и посмотрим как будет меняться Пуассоновское поле и уровень капитала
\begin{center}
\includegraphics[width=0.7\linewidth]{мой отчёт/44.png}
\end{center}
#
В ходе работы были реализованы модели двумерных пуассоновских полей для различных интерпретаций. Показано, что поведение системы существенно зависит от параметров.
\begin{thebibliography}{10}

\bibitem{smirnov} 
Смирнов С.\,Н. Курс лекций <<Стохастический анализ и моделирование>>, 2025.

\bibitem{ulyanov} 
Ульянов В.\,В. Курс лекций <<Теория вероятностей и математическая статистика>>, 2023--2024.

\bibitem{kropacheva} 
Кропачева Н.\,Ю., Тихомирова А.\,С. <<Моделирование случайных величин: Метод. указания>>. НовГУ им. Ярослава Мудрого, 2004.

\bibitem{bulinski} 
Булинский А.\,В., Ширяев А.\,Н. <<Теория случайных процессов>>. М.: ФИЗМАТЛИТ, 2005.

\bibitem{sobol} 
Соболь И.\,М. <<Численные методы Монте-Карло>>. Издание <<Наука>>, главная редакция физико-математической литературы, 1973.

\bibitem{ventzel} 
Вентцель А.\,Д. <<Курс теории случайных процессов>>, 2-е изд., доп. --- М.: Наука. Физматлит, 1996.

\bibitem{gasnikov} 
Гасников А.\,В. <<Лекции по случайным процессам>>. Москва, МФТИ, 2019.

\bibitem{feller} 
Феллер В. <<Введение в теорию вероятностей и её приложения>>. Т.1. М., Мир, 1984.

\bibitem{digaylova} 
Дигайлова И.\,А. <<Теория идентификации>>, 2012.

\bibitem{lewis} 
P.\,A.\,W. Lewis, G.\,S. Shedler. ``Simulation of nonhomogeneous poisson processes by thinning''. IBM Research Laboratory, San Jose, California, 1979.

\end{thebibliography}

\end{document}